\section{Methodology}
\label{sec:methodology}

\subsection{Research questions}
This paper asks two questions. \textbf{RQ1:} using only black-box,
coverage-free feedback, does LLM-guided refinement of a grammar-based
generator improve that generator's acceptance rate and structural diversity
relative to a static baseline, on a fixed target? \textbf{RQ2:} does the
pipeline -- harness contract, proxy signal, refinement loop, and sandbox --
generalize to a second, independently adapted target without modification?

\subsection{Target, grammar, and harness}
The primary target is cJSON, pinned at v1.7.19
(commit \texttt{c859b25}). The grammar is the ANTLR grammars-v4 JSON
grammar~\cite{antlr2024grammarsv4}, vendored unmodified. A thin C harness
reads a candidate input from stdin and parses it with
\texttt{cJSON\_ParseWithLengthOpts(..., require\_null\_terminated=1)}, the
library's strictest entry point: it rejects any input with trailing bytes
after a valid value, so an ``accepted'' input is one where the entire buffer
is a single grammar-valid value, matching the grammar's
\texttt{json : value EOF ;} production exactly. The library's more commonly
used \texttt{cJSON\_Parse} entry point does not enforce full-buffer
consumption and was deliberately not used, since it would silently accept
inputs the grammar does not.

Reconciling the formal grammar with the target's actual behavior surfaced
three adaptations, each a genuine mismatch between what the grammar
specifies and what the implementation does: (i) the grammar has no token for
\texttt{NaN}/\texttt{Infinity}, matching cJSON's inability to produce or
accept them, so generators exclude non-finite floats; (ii) cJSON accepts
duplicate object keys, storing every entry in an internal linked list rather
than rejecting or de-duplicating, even though the grammar is silent on
duplicate keys -- this is treated as valid-but-underspecified territory, not
a rejection case, and generators deliberately exercise it; (iii) the
full-buffer-consumption requirement above is enforced by the harness's
choice of entry point, not by the grammar itself. Section~\ref{sec:evaluation}
reports the analogous set of adaptations found for the second target,
parson, where two of the three go the opposite direction under the identical
grammar.

\subsection{Outcome classification}
Every input is executed as a bounded subprocess (5-second default timeout)
against a sanitizer-instrumented build (\texttt{-fsanitize=address,undefined}).
The result is classified, in order: \texttt{timeout} if the process does not
return within the bound; \texttt{signal:<name>} if it is killed by a fatal
signal; \texttt{crash} if it exits with a non-zero return code or its
stderr contains \texttt{AddressSanitizer} or \texttt{UndefinedBehaviorSanitizer}
text; \texttt{accepted} if its stdout begins with \texttt{status=accepted};
otherwise \texttt{rejected}. Timeouts are classified and reported alongside
crashes rather than discarded, since an unbounded hang in a parser is a
denial-of-service bug in its own right.

\subsection{A coverage-free proxy signal}
With no coverage instrumentation, refinement is steered by three signals
computed purely from a campaign's classified outcomes:

\begin{itemize}
  \item \textbf{Acceptance rate} -- the fraction of inputs classified
    \texttt{accepted}. A generator accepting close to 0\% of its own output
    is testing only the rejection path, not the parser's structural
    handling of valid input, and is flagged for correction.
  \item \textbf{Structural diversity} -- the count of distinct
    \emph{structural fingerprints} seen in a campaign. A fingerprint maps
    each byte of an input to one of six symbols (the six structural
    characters \texttt{\{\}[]:,} are kept literal, whitespace maps to
    \texttt{\_}, ASCII digits map to \texttt{\#}, the double-quote maps to
    \texttt{"}, and every other byte maps to \texttt{x}), concatenates the
    result, and truncates it to 256 characters. Two inputs share a
    fingerprint exactly when their first 256 bytes reduce to the same
    positional shape string; the count of distinct fingerprints across a
    campaign is used as a coverage-free proxy for how many different
    grammar shapes a generator has actually produced. The 256-byte
    truncation makes this proxy a lower bound on true shape diversity for
    inputs longer than that, since two structurally different inputs that
    agree on their first 256 bytes collapse to one fingerprint; this is a
    deliberate accuracy/cost tradeoff and is stated as a limitation in
    Section~\ref{sec:discussion}.
  \item \textbf{Rejection signatures} -- the count of distinct parser
    rejection messages and offsets seen, used to tell the proposer which
    malformed shapes are already well covered versus unexplored.
\end{itemize}

These three numbers, plus outcome counts and unique input lengths, are the
entire feedback surface exposed to the LLM each iteration; no coverage data,
execution trace, or source excerpt is ever included in the prompt.

\subsection{Refinement loop}
Each iteration builds a prompt containing the grammar text and a summary of
the previous iteration's campaign (outcome counts, unique input lengths,
unique structural fingerprints, unique rejection signatures), and sends it
to a proposer -- an LLM in the live configuration, or, for the parson
generalization test, a human reasoning over the same summary and playing
the identical role (Section~\ref{sec:evaluation}). The proposer must return
a single Python function, \texttt{generated\_json}, decorated with
Hypothesis's \texttt{@st.composite} and built from \texttt{st.*}
combinators, so that recursive structure is expressed the same way a
human-authored Hypothesis strategy would express it. The official run uses
\texttt{gpt-4.1-mini} at temperature 0.2 as the proposer.

\subsection{Sandboxing untrusted, LLM-authored code}
Because the proposer's output is executable Python that this pipeline runs
locally, every proposal passes through a validator
(\texttt{load\_strategy}) before it ever touches the target binary. The
validator first parses the proposal's AST and rejects it outright if it
imports anything other than \texttt{hypothesis.strategies}, or if it calls
\texttt{eval}, \texttt{exec}, \texttt{open}, \texttt{compile}, or
\texttt{\_\_import\_\_} by name. The proposal is then executed with an
otherwise ordinary Python builtin namespace, minus a short, explicit
blocklist (\texttt{eval}, \texttt{exec}, \texttt{compile},
\texttt{\_\_import\_\_}, \texttt{open}, \texttt{input}, \texttt{breakpoint},
\texttt{exit}, \texttt{quit}, \texttt{help}, \texttt{globals},
\texttt{locals}, \texttt{vars}, \texttt{getattr}, \texttt{setattr},
\texttt{delattr}) -- ordinary helpers such as \texttt{ord}, \texttt{chr},
\texttt{str}, \texttt{len}, and \texttt{range} remain available, so a
proposal is free to use normal Python rather than being restricted to
Hypothesis combinators alone. Finally, the resulting \texttt{generated\_json}
must be callable and must produce a \texttt{bytes} value via
\texttt{.example()}; a proposal that fails any of these checks is rejected
and logged, and never reaches the harness. This is deliberately scoped as a
filter against \emph{accidental} misuse of untrusted, LLM-authored code, not
as a hardened security boundary: Python's object model can reach
functionality beyond any blocklisted name (for instance, by walking
\texttt{().\_\_class\_\_.\_\_base\_\_.\_\_subclasses\_\_()}), and closing
that off would require process- or container-level isolation rather than
name filtering. We report this limitation directly rather than implying a
guarantee the design does not provide.

\subsection{Crash triage and minimization}
A sanitizer report is reduced to its \texttt{SUMMARY:} line and stack
frames with addresses and source line numbers normalized out, then hashed
into a 16-character signature; records sharing a signature are treated as
one root cause. For each unique signature, Hypothesis's \texttt{find} is
used to shrink a triggering input to a minimal byte string that still
reproduces the same crash or timeout classification -- rather than simply
keeping the first crash observed -- and the minimized input is re-run once,
standalone, against the pinned build to verify it reproduces outside the
shrinking process itself.

\subsection{Reproducibility engineering}
Hypothesis's \texttt{.example()} draws from a private internal random state
that Python's own \texttt{random.seed()} does not reach, and separately
orders its candidate batch with the module-level \texttt{random.shuffle}.
Seeding only one of the two leaves the produced example stream
non-reproducible across process invocations; a shim
(\texttt{seeding.seed\_everything}) seeds both, which we verified gives
byte-identical example streams across repeated processes on Hypothesis
6.165.5 / CPython 3.14.7. Two documented, non-internal alternatives --
\texttt{settings(derandomize=True)} and \texttt{global\_force\_seed} -- were
rejected because both collapse a run's draws into repeats of a single small
batch, which would change the generation algorithm itself rather than only
fix its starting point. Because this shim depends on an internal
implementation detail, the Hypothesis version is pinned in
\texttt{requirements.txt} and the shim fails loudly rather than silently if
the relied-upon attribute is renamed or removed in a future release. Every
run additionally persists a manifest (seed, package versions, platform,
model, feedback mode, CLI arguments) and a reproducibility report (compiler,
harness SHA-256, exact command line, repository commit); the compiled
harness binary itself is gitignored and host-compiled, so it is identified
by content hash rather than assumed identical across machines. This
reproducibility guarantee is scoped to the generated input stream only: it
does not extend to LLM proposals themselves (temperature 0.2, no server-side
seed), to wall-clock-dependent timeout classification, or to the harness
binary's bit-for-bit identity across hosts and toolchains.

\subsection{Experimental design}
For RQ1, we run five independent, seeded trials (seeds 0--4) of a static
baseline generator and five independent, seeded trials of the LLM-refined
loop against the same pinned cJSON build, each trial consisting of up to
five refinement iterations of up to 500 examples. We report acceptance
rate, structural fingerprints, and rejection signatures per trial, and
compare the two arms with a non-parametric two-sample test appropriate for
five observations per arm without assuming a particular distribution
(Section~\ref{sec:evaluation}); we do not assume, and do not claim, that
count-based metrics are directly comparable between arms, since the two
arms' trials execute different total numbers of inputs. For RQ2, we run the
identical pipeline, entirely unmodified, against parson for five real
iterations, with a human reasoning over each iteration's campaign summary
in the proposer's role and every resulting proposal passing through the
same \texttt{load\_strategy} validation and \texttt{run\_campaign} execution
path a model's output would take; this substitutes for a live API call
without altering any other part of the pipeline being tested.
